\subsection{Contratos y firma bilateral}
\label{sec:impl-firma-bilateral}

Esta subsección implementa el diseño de la \cref{sec:diseno-firma-bilateral}
para CU-13, CU-14 y CU-18 (\cref{tab:cu-13,tab:cu-14,tab:cu-18}) y para los
criterios bilaterales de RF-15 a RF-17 (\cref{tab:rf-cripto}). El estado que
describe es \emph{implementado y demostrado localmente}: el servidor, el
protocolo, la evidencia, el validador, las migraciones, los avisos y la
demostración reproducible existen y se ejecutaron sobre servicios desechables;
ningún paso se ha instalado en producción ni se ha ensayado con un Cliente
real, por lo que no se presenta como aceptado.

\paragraph{Partes y alta del Cliente.} La migración 0050 sustituye las dos
funciones de vínculo de firma documental para que un Owner, un Litigante o un
Cliente activos puedan registrar su certificado; Paralegal sigue excluido y no
cambia ninguna tabla ni fila. La firma individual conserva su exigencia de un
Litigante miembro del expediente, de modo que un vínculo de Owner o de Cliente
nunca produce una firma individual. La migración
\texttt{0051\_client\_invitations} amplía la restricción de rol de las
invitaciones para admitir \texttt{client}, solo emitidas por un Owner, y el
arranque exige la restricción ampliada. Solo una cuenta creada por una
invitación \texttt{client} aceptada es elegible como Cliente de una solicitud:
una cuenta de alta directa no lo es, porque el Owner fijó su contraseña
inicial y su secreto TOTP se generó en el servidor. La migración
\texttt{0051\_client\_memberships} añade el disparador que impide asignar una
cuenta Cliente a un expediente; el servicio de expedientes responde
\texttt{422 invalid\_input}, los selectores de miembros omiten esas cuentas y
el cambio de rol de una cuenta asignada hacia Cliente se rechaza. La migración
se detiene si encuentra una membresía de Cliente y nunca la borra.

\paragraph{Almacenamiento.} La familia de contratos se divide en tres archivos
solo para respetar el límite de 400 líneas: \texttt{0052\_contracts} (tablas y
constructores de las declaraciones), \texttt{0052\_contract\_events} (bloqueo,
preservación, eventos y solicitudes) y \texttt{0052\_contract\_signatures}
(firmas, compleción, emparejamiento, avisos, guarda de versión y disparadores).
Las seis tablas de la \cref{sec:diseno-contratos-datos} no tienen columna de
estado; los disparadores reconstruyen en SQL los 240 bytes de
\texttt{CTRTREQ1} y los 303 de \texttt{CTRTSIG1}, exigen el turno, la vigencia,
la cuenta y el vínculo designados y la misma raíz en ambas posiciones, y
toman el bloqueo consultivo de los vínculos después del bloqueo auditado.
La guarda de versión rechaza el sello y la firma individual sobre una versión
con solicitud abierta o completada, y la firma rechaza una versión que ya tenga
evidencia, de modo que \texttt{documents.evidence} permanece nulo en toda
versión contratada. Las referencias nuevas a la auditoría elevan sus
disparadores de clave foránea de 26 a 28. Al arrancar, el inventario comprueba
en SQL cada predicado, revisión, evento y byte reconstruido; después decodifica
cada declaración con los códecs del dominio, verifica cada firma con la
revisión de confianza que nombra y los tres sellos sobre sus huellas, y
compara cada manifiesto con lo almacenado.

\paragraph{Servicio y rutas.} El módulo de aplicación de contratos deriva el
estado y decide cada acceso a partir de los hechos que entrega el almacén, con
la definición única de solicitud abierta. La API expone catorce rutas bajo
\texttt{/api/v1}: seis por expediente para listar, crear, consultar, solicitar,
abandonar y recuperar el recibo de creación o abandono, y ocho por solicitud
para la bandeja propia, el detalle, el PDF exacto, la preparación sin
escritura, la firma, el cierre por rechazo, declinación o cancelación, el
recibo de una operación y los paquetes parcial y completo. El total de rutas
pasa de 223 a 237 y la matriz de acceso registra cada par. Una cuenta sin
relación recibe \texttt{404 contract\_not\_found}; una cuenta relacionada que
no es la actora de la acción, \texttt{403}. Toda mutación lleva
\texttt{operation\_id} y revisión esperada, devuelve un recibo exacto y no se
reintenta de forma automática; la repetición idéntica devuelve el recibo
original. La TSA se consulta fuera de la transacción, una vez por envío: T1
sobre la firma del Cliente y, en la contrafirma, Td sobre el PDF y Tf sobre el
manifiesto, con \texttt{T1 <= Td <= Tf} y la vigencia comprobada con
\texttt{genTime} más su exactitud. Una falla de la TSA responde
\texttt{503 contract\_timestamp\_unavailable} sin escribir estado. El
certificado del Cliente revocado en la CRL vigente bloquea la contrafirma
(\texttt{contract\_counterparty\_revoked}); el retiro posterior de su vínculo,
no. Las listas, detalles e historiales documentales devuelven el estado de
contrato derivado, el filtro de pendientes de sello omite las versiones
congeladas y el indicador del tablero cuenta contratos en borrador o en firma.

{\emergencystretch=3em \paragraph{Evidencia.} La exportación arma el paquete parcial
\texttt{BilateralContractPartialV1}, de 14 entradas, disponible para el
Cliente desde T1 y también tras un cierre sin éxito, y el completo
\texttt{BilateralContractV1}, de 26 entradas, ambos con
\texttt{INSTRUCCIONES.md}. El manifiesto de la contrafirma se codifica con los
bytes exactos que escribe la exportación, y esta se niega a entregar un
paquete cuyos artefactos difieran del manifiesto guardado; exportar dos veces
produce los mismos bytes. El servidor usa los códecs de
\texttt{evidence-verify}, el mismo motor del validador
(\cref{sec:impl-validador}), y las instrucciones reproducen con OpenSSL cada
firma, huella y sello, incluida la lectura de la huella del PDF en el
desplazamiento 206 de cada declaración.\par}

\paragraph{Vencimiento y avisos.} El servidor compone siempre un consumidor
supervisado de contratos. Cada cinco segundos escribe, con el actor técnico
\texttt{system:contract-expiry}, el evento de vencimiento de hasta 20
solicitudes agotadas y sus avisos. Los avisos nombran el título, la solicitud y
el vencimiento, sin token, firma ni certificado. Su envío por correo es un
consumidor arrendado aparte que solo actúa con
\texttt{TT\_CONTRACT\_EMAIL\_ENABLED=true} y reutiliza la configuración del
correo de informes; deshabilitado, los avisos quedan guardados y las partes
ven sus solicitudes en Qadra. Nada del flujo de firma espera al correo.

\paragraph{Interfaz por rol.} En Qadra, el Owner y el Litigante ven la acción
\emph{Crear contrato} en la ficha de un documento PDF sin sellar y sin
contrato, y la vista \emph{Contratos} de cada expediente; un asistente sobre
la versión vigente fija el Cliente por su correo, el firmante del despacho, la
vigencia de siete días por omisión y el folio de la hoja de huella, y muestra
la lista de preparación que devuelve el servidor antes de solicitar firmas. El detalle del contrato
presenta su estado, la línea de tiempo, las solicitudes con sus descargas y el
aviso de una versión posterior; el solicitante o un Owner cancelan con motivo
obligatorio y confirmación aparte. Un Paralegal asignado solo lee ese detalle.
La lista, la ficha y el historial documentales leen una versión sin sellar de
una serie en firma como \emph{En contrato} o \emph{Contrato firmado}, nunca como
sellada o pendiente de sello, y ocultan el sello y la firma individual. El
Cliente solo dispone de \emph{Mis contratos}, \emph{Mi firma} y
\emph{Validar evidencia}; cualquier otra ubicación lo devuelve a sus contratos.
Firma en tres pasos: revisa el PDF, cuya huella el navegador recalcula y
compara con la solicitud firmada antes de mostrarlo, junto con la contraparte,
su huella de certificado y el vencimiento; confirma una casilla en lenguaje
llano, y firma con su archivo de clave y contraseña sobre los 303 bytes
preparados, sin que la clave salga del navegador. Rechazar es una acción
separada con confirmación. El firmante del despacho encuentra la solicitud en
\emph{Por firmar}, revisa el mismo PDF, descarga el acuse parcial y lo valida
en el navegador; la contrafirma permanece deshabilitada hasta que el informe
local corresponde al perfil parcial, no presenta fallas y cubre los bytes
exactos de la solicitud. Antes de firmar, el navegador comprueba además que la
declaración preparada encadena con la firma del Cliente y con T1 tal como el
servidor los entrega, bajo el mismo despliegue y la misma raíz. Declinar sigue
disponible sin validar. En \emph{Mi firma}, la generación de clave muestra la
huella SHA-256 de la clave pública calculada en el navegador y ofrece una hoja
de huella para imprimir con el correo de la cuenta y la huella de la raíz. Una
respuesta incierta se concilia con el recibo de su operación mientras la
página sigue abierta, y una declaración rechazada por cambio de cuenta, de
confianza o de revisión se prepara de nuevo. El recorrido completo se
reproduce además por API con \rutarepo{scripts/api-contract-demo.sh}
(\cref{sec:pruebas-firma-bilateral}).

\paragraph{Límites.} La firma es electrónica simple con TSA local RFC~3161; no
hay NOM-151, PSC, anclaje externo, FIREL ni PAdES. Un contrato admite
exactamente dos partes en orden fijo. El servidor nunca recibe la clave privada
de ninguna parte. Los límites (a) a (i) de la \cref{sec:diseno-firma-bilateral}
se aplican sin excepción. El recorrido en el servicio instalado se describe en
la \cref{sec:pruebas-firma-bilateral}; la ceremonia con un cliente real y la
evaluación con personas siguen pendientes.

\Fuentes{ADR 0086, \rutarepo{docs/http-contract-signing.md},
\rutarepo{crates/application/src/contracts/},
\rutarepo{crates/infrastructure/src/contract_schema/} y
\rutarepo{migrations/0052_contracts.sql}.}
